
This work compared matrix-level and token-level distillation objectives for Transformer-to-SSM conversion, measuring hidden state drift across sequence lengths 512--2048 tokens using a unified Phi-Mamba framework.

Key findings:
\begin{itemize}
\item CAB (token-level) drift slope: 0.00090506
\item MOHAWK (matrix-level) drift slope: 0.00452495
\item Ratio: $5\times$
\item CAB drift ratio bounded at 1.34; MOHAWK at 2.02
\item Slope difference statistically significant (p < 0.001)
\end{itemize}

These measurements were obtained using simulated student models. The drift pattern is consistent with the hypothesis that token-level supervision is position-agnostic while matrix-level supervision encodes length-dependent patterns. Downstream F1 retention was not verified.

Phi-1.5's 2048-token hard limit constrains the experimental scope. Length extrapolation analysis at 16K--32K requires teachers with relative position encodings.

\subsubsection{Future Work}

\textbf{Full training verification}: Full 1.5B token training per condition would test whether the drift stability gap translates to task performance at extrapolated lengths.

\textbf{RoPE-based teachers}: Llama, Mistral, and other RoPE-based models can process longer sequences natively, enabling true length extrapolation experiments.

\textbf{Real student checkpoints}: Validation with actual trained phi-mamba and CAB-distilled checkpoints would confirm the drift magnitudes observed in simulation.

